% !TEX root = main.tex

\documentclass[../main.tex]{subfiles}
\begin{document}

\section{Sicco}
\label{sec:sicco}

\begin{povolani}[Sicco]
\textbf{Svolení}:\\
\textbf{Cena}:\\
\textbf{Popis}: Speciální dovednosti (\dov{Etiketa}, \dov{Falšování}) a speciality (\trk{Znalost ulice}, \trk{Sítě}, \trk{Lež za lží}, \trk{Zoon politikon}, \trk{Reputace}, \trk{Dvojí identita}), viz níže.
\end{povolani}

\subsection{Speciální dovednosti}
\label{sec:sicco-dovednosti}

\subsubsection{Etiketa}
\label{sec:sicco-etiketa}
\dov{Etiketa} (I) – zahrnuje společenské chování.

\subsubsection{Falšování}
\label{sec:sicco-falsovani}
Sicco se naučí psát a vyjadřovat jako někdo jiný. Dokáže tak efektivně napodobovat dopisy, pečetě, potvrzení, smlouvy atd.

\begin{akce}
  \iconitem{prekonani-ikona}{\akc{Překonání:} Kdykoliv, kdy Sicco chce někoho oblafnout falzifikátem, používá \dov{Falšování} jako \akc{Překonání}.}
  \iconitem{vytvoreni-ikona}{\akc{Vytvoření výhody:} Při úspěšném zfalšování se na oblafnutých jedincích umístí aspekt typu \asp{Důvěřivý špipík}, \asp{Neví, která bije} atd. Podobně lze takto umisťovat aspekty na předměty – učinit dopis \asp{Zfalšovaným} atd.}
  \iconitem{utok-ikona}{\akc{Útok:} Tuto dovednost nelze použít k vedení útoku.}
  \iconitem{obrana-ikona}{\akc{Obrana:} Tuto dovednost nelze použít k obraně.}
\end{akce}


\subsection{Speciality}
\label{sec:sicco-speciality}

\begin{specialita}[Znalost ulice, II]
\begin{itemize}
	\item \textbf{Popis:} Sicco je v ulicích města jako doma. Při konfliktech, které se odehrávají v městském prostředí tak z něj přirozeně čerpá výhody – v každém konfliktu lze jednou zdarma vyvolat aspekt odpovídající prostředí (\asp{Oči na každém rohu}, \asp{Temné uličky}, \asp{Není se kde schovat}).
\end{itemize}
\end{specialita}

\begin{specialita}[Sítě, III]
\begin{itemize}
	\item \textbf{Popis:} Sicco buduje síť svých kontaktů. Může si tak udržovat přátele blízko a nepřátele ještě blíž. Svojí Síť si může libovolně pojmenovat, postavy do ní lze zaplést jednorázově hodem na relevantní dovednost (sociální) s úspěchem se stylem. K tomu lze přirozeně používat odhalené aspekty na kartě dotyčných postav stejně jako jiné speciality – \trk{Lež za lží}, \trk{Reputace} atd.
	\item \textbf{Vylepšení sítě:} za každý dovednostní bod lze dokoupit další Vylepšení sítě:
		\begin{itemize}
			\item \trk{Poslíčci}: Postavy propojené Síti mezi sebou dokáží komunikovat mnohem rychleji; využívají služeb profesionálních Poslíčků. Ve městech mají spolehlivé Poslíčky, kteří jsou jim k ruce, dokáží spolehlivě předávat zprávy mimo městy.
			\item \trk{Společenský polštář}: Sicco může členům své Sítě poskytnout jedno políčko duševního stresu o hodnotě 2 a umožnit jim tak lépe ustát společenské či duševní útoky.
			\item \trk{Pod dohledem}: Sicco si udržuje povědomí o aktivitách a stavu postav zapletených v jeho Síti aniž by musel vyvinout větší iniciativu, má přístup k základním informacím ohledně postav, tj. např. zdravotní stav, místo pobytu apod. Ty pak mohou sloužit jako pobídka k dalšímu zjišťování.
			\item \trk{Zdroje}: Sicco získává 3 políčka „majetkového stresu“, který může využít k financování svých aktivit. Při použití se políčko zaškrtne a obnoví se až po určité době. Sicco může použít své vlastní peníze jako okamžitou finanční úlevu a ihned si políčko „vyléčit“.
				\begin{itemize}
					\item \textbf{Službička} – obnoví se po každém sezení; zahrnuje služby cca v rozsahu jednotek zl.
					\item \textbf{Práce} – obnoví se jednou za dobrodružství; zahrnuje služby cca v rozsahu desítek zl.
					\item \textbf{Peníze hýbou světem} – obnoví se jednou za příběhový oblouk (cca); zahrnuje služby cca v rozsahu stovky zl.
				\end{itemize}
		\end{itemize}
\end{itemize}
\end{specialita}

\begin{specialita}[Lež za lží, III]
\begin{itemize}
	\item \textbf{Popis:} Při vytváření výhody \dov{Klamáním} s úspěchem alespoň +2 se automaticky na cíl navíc umístí aspekt \asp{V síti lží}. Při jakémkoliv dalším pokusu o \dov{Klam} (ať už akcí \akc{Překonání} či \akc{Vytvoření výhody}) získává Sicco na tuto postavu jedno volné vyvolání na tento aspekt; lze použít jednou za sezení.
\end{itemize}
\end{specialita}

\begin{specialita}[Zoon politikon, IV]
\begin{itemize}
	\item \textbf{Popis:} Skrze působení ve společnosti a zapletení do společenských konfliktů získává Sicco jednak jistou odolnost – může si připsat jedno políčko duševního stresu o hodnotě 2, druhak získává pro účely duševního zranění (např. skrze dovednost \dov{Provokace}) hodnocení zbraně +1.
\end{itemize}
\end{specialita}

\begin{specialita}[Reputace, IV]
\begin{itemize}
	\item \textbf{Popis:} Sicco získá na svoji kartu nový aspekt postavy – se vším všudy. Ten odráží jeho Reputaci, jak ho vnímá okolí. Přirozeně, toto je zcela v siccových rukou, může si vybásnit ledasco. Toto lze prokouknout hody na \dov{Empatii}, \dov{Vztahy}, \dov{Kontakty} atd.
\end{itemize}
\end{specialita}

\begin{specialita}[Dvojí identita, V]
\begin{itemize}
	\item \textbf{Popis:} Přístupem k mocnému magickému předmětu Sicco získává možnost zcela změnit svoji totožnost – rasa, vzhled, hlas a vše další. To dokáže udělat okamžitě. Všechny aspekty, které Sicco takto umístí na jiné postavy či objeví jsou přístupné oběma jeho identitám; naopak, aspekty umístěné na něj jsou spojeny s identitou, kterou měl zrovna sicco na sobě. Přirozeně, vlastní-li Sicco specialitu \trk{Reputace}, smí si pod druhou identitou změnit i tento aspekt. Pro jakékoliv mechanické potřeby lze Druhou identitu považovat za zcela nezávislou postavu.
\end{itemize}
\end{specialita}


\end{document}

%%% Local Variables:
%%% mode: LaTeX
%%% TeX-master: "../../main"
%%% End:
